% ==========================================================
% DOCUMENTO TÉCNICO - PROYECTO INTEGRADOR
% Universidad Estatal Península de Santa Elena
% Sistema Web de Gestión de Incidencias Georreferenciadas
% ==========================================================

\documentclass[12pt,a4paper]{article}

% ==========================================================
% PAQUETES
% ==========================================================
\usepackage[utf8]{inputenc}
\usepackage[T1]{fontenc}
\usepackage[spanish,es-tabla]{babel}
\usepackage[top=2.25cm,bottom=2.25cm,left=2.25cm,right=2.25cm]{geometry}
\usepackage{graphicx}
\usepackage[table]{xcolor}
\usepackage{titlesec}
\usepackage{enumitem}
\usepackage{tabularx}
\usepackage{booktabs}
\usepackage{longtable}
\usepackage{array}
\usepackage{float}
\usepackage{hyperref}
\usepackage{parskip}
\usepackage{fancyhdr}
\usepackage{lastpage}
\usepackage{microtype}
\usepackage{amsmath}
\usepackage[most]{tcolorbox}
\usepackage{pdflscape}

% ==========================================================
% COLORES Y ESTILO
% ==========================================================
\definecolor{primarygray}{HTML}{4A4A4A}
\definecolor{secondarygray}{HTML}{E8E8E8}
\definecolor{softgray}{HTML}{F3F6F8}
\definecolor{textgray}{HTML}{2B2B2B}

\hypersetup{
    colorlinks=true,
    linkcolor=primarygray,
    urlcolor=primarygray,
    citecolor=primarygray
}

\titleformat{\section}
{\color{primarygray}\Large\bfseries}
{\thesection.}
{0.55em}
{}
[\titlerule]

\titleformat{\subsection}
{\color{primarygray}\large\bfseries}
{\thesubsection.}
{0.45em}
{}

\pagestyle{fancy}
\setlength{\headheight}{14pt}
\fancyhf{}
\lhead{\small Documento Técnico -- Proyecto Integrador}
\rhead{\small Sistema SGI}
\cfoot{\small Página \thepage\ de \pageref{LastPage}}
\renewcommand{\headrulewidth}{0.4pt}
\renewcommand{\footrulewidth}{0pt}

\renewcommand{\arraystretch}{1.12}
\setlength{\tabcolsep}{4.2pt}
\setlist[itemize]{leftmargin=*, itemsep=1pt, topsep=2pt}
\setlist[enumerate]{leftmargin=*, itemsep=1pt, topsep=2pt}
\setlength{\parskip}{4pt}
\sloppy

\setcounter{tocdepth}{2}
\setcounter{secnumdepth}{2}

% ==========================================================
% COMANDO PARA EVIDENCIAS INDIVIDUALES
% ==========================================================
\newcommand{\evidencia}[4]{
\clearpage
\begin{figure}[H]
    \centering
    \IfFileExists{#1}{
        \includegraphics[width=#2\textwidth,height=#3\textheight,keepaspectratio]{#1}
    }{
        \begin{tcolorbox}[
            colback=secondarygray,
            colframe=primarygray,
            boxrule=0.8pt,
            width=0.90\textwidth,
            height=5cm,
            halign=center,
            valign=center
        ]
            \textit{[Insertar archivo: \texttt{#1}]}
        \end{tcolorbox}
    }
    \caption{#4}
\end{figure}
}

% ==========================================================
% COMANDO PARA IMÁGENES CON RESPALDO
% ==========================================================
\newcommand{\captura}[3]{
\begin{figure}[H]
    \centering
    \IfFileExists{#1}{
        \includegraphics[width=#2\textwidth,keepaspectratio]{#1}
    }{
        \begin{tcolorbox}[
            colback=secondarygray,
            colframe=primarygray,
            boxrule=0.8pt,
            width=#2\textwidth,
            height=3.6cm,
            halign=center,
            valign=center
        ]
            \textit{[Insertar archivo: \texttt{#1}]}
        \end{tcolorbox}
    }
    \caption{#3}
\end{figure}
}


% ==========================================================
% DOS FIGURAS INDEPENDIENTES EN UNA MISMA PÁGINA
% ==========================================================
\newcommand{\dosporevidencia}[6]{
\clearpage
\begin{figure}[H]
    \centering
    \includegraphics[width=#2\textwidth,height=0.34\textheight,keepaspectratio]{#1}
    \caption{#3}
\end{figure}

\vspace{0.15cm}

\begin{figure}[H]
    \centering
    \includegraphics[width=#5\textwidth,height=0.34\textheight,keepaspectratio]{#4}
    \caption{#6}
\end{figure}
}

\begin{document}

% ==========================================================
% PORTADA -- CONSERVADA COMO FUE ENTREGADA
% ==========================================================
\begin{titlepage}
\centering
\vspace*{1.0cm}

% --- Encabezado institucional ---
\IfFileExists{upse.jpg}{
    \includegraphics[width=3.8cm]{upse.jpg}\\[0.6cm]
}{
    {\Large \textbf{UPSE}}\\[0.6cm]
}

{\large \textbf{\textcolor{primarygray}{UNIVERSIDAD ESTATAL PENÍNSULA DE SANTA ELENA}}}\\[0.18cm]
{\normalsize \textbf{Facultad de Sistemas y Telecomunicaciones}}\\[0.06cm]
{\normalsize \textbf{Carrera de Software}}\\[0.06cm]
{\normalsize Sexto Semestre}

\vspace{0.9cm}
\noindent\rule{\textwidth}{0.8pt}
\vspace{0.3cm}

% --- Asignatura ---
{\normalsize \textsc{Asignatura:} \textbf{Tecnología y Desarrollo Web}}

\vspace{0.3cm}
\noindent\rule{\textwidth}{0.8pt}
\vspace{0.8cm}

% --- Tipo de documento ---
{\large \textbf{\textcolor{primarygray}{DOCUMENTO TÉCNICO DEL PROYECTO INTEGRADOR}}}

\vspace{0.8cm}

% --- Título del proyecto ---
{\LARGE \textbf{Sistema Web de Gestión de}}\\[0.2cm]
{\LARGE \textbf{Incidencias Georreferenciadas}}\\[0.3cm]
{\large \textit{(SGI)}}

\vfill

% --- Datos académicos ---
\begin{tabular}{@{}r@{\quad}l@{}}
\textbf{Integrantes:}          & Damián Jesús Torres Cadena \\
                                & Melanie Adriana Tomalá Merejildo \\
                                & Jean Pierre Cedeño Andrade \\[0.4cm]
\textbf{Docente:}              & Ing. Evelyn Del Pezo Izaguirre \\[0.4cm]
\textbf{Repositorio:}          & \small\url{https://github.com/DevTorres404/georeferenced-incident-management} \\[0.15cm]
\textbf{Sistema en producción:}& \small\url{https://app.labtorres.me/} \\[0.4cm]
\textbf{Período académico:}    & 2026--1 \\
\textbf{Fecha de entrega:}     & Julio de 2026 \\
\end{tabular}

\vspace{1.0cm}

\end{titlepage}

\clearpage
\pagenumbering{arabic}
\setcounter{page}{1}

% ==========================================================
% ÍNDICE GENERAL
% ==========================================================
\tableofcontents
\clearpage

\listoffigures
\clearpage

\listoftables
\clearpage

% ==========================================================
% INFORMACIÓN ACADÉMICA COMPLEMENTARIA
% ==========================================================
\section*{Información académica e integración del proyecto}
\addcontentsline{toc}{section}{Información académica e integración del proyecto}

El proyecto integrador articula los resultados de aprendizaje de las asignaturas
\textbf{Tecnologías y Desarrollo Web}, \textbf{Calidad de Software},
\textbf{Administración de Data Center} y \textbf{Base de Datos I y II}. La solución
presentada integra frontend, backend, persistencia relacional y geoespacial,
despliegue mediante contenedores y evidencias de validación y calidad.

El trabajo se desarrolló de forma colaborativa, manteniendo la integración entre
la interfaz, los servicios de negocio, la base de datos y la infraestructura. Las
secciones siguientes describen la implementación realizada, las decisiones
adoptadas, la forma de ejecución y las evidencias obtenidas.

% ==========================================================
% 1. DESCRIPCIÓN DE LA IMPLEMENTACIÓN
% ==========================================================
\section{Descripción de la implementación}

El \textbf{Sistema Web de Gestión de Incidencias Georreferenciadas (SGI)} se desarrolló como una aplicación web completa para registrar, consultar, asignar y dar seguimiento a incidencias asociadas a una ubicación geográfica. La solución integra un frontend web, una API REST, persistencia relacional con capacidades geoespaciales y un entorno de producción basado en contenedores.

En el backend se utilizó una organización por dominios con capas de \textbf{Dominio, Aplicación e Infraestructura}, siguiendo principios de Arquitectura Hexagonal y DDD. Esta separación mantiene reglas como estados, clasificación, asignaciones, permisos y reaperturas fuera de los detalles de HTTP y persistencia. El frontend se organizó mediante módulos JavaScript ES6 y centraliza el consumo de la API en un cliente HTTP común.

Las decisiones de implementación más relevantes fueron:
\begin{itemize}
    \item \textbf{Frontend y backend desacoplados:} la interfaz se publica de forma independiente y consume la API mediante HTTP/JSON.
    \item \textbf{Trazabilidad:} se conserva historial de estados, asignaciones, clasificación, comentarios y ciclos de resolución.
    \item \textbf{Georreferenciación:} las incidencias almacenan latitud, longitud y geometría \texttt{Point} mediante PostGIS.
    \item \textbf{Procesamiento asíncrono y tiempo real:} Redis soporta caché, sesiones y colas; Laravel Reverb distribuye eventos mediante WebSockets.
    \item \textbf{Evidencias digitales:} los adjuntos se almacenan mediante RustFS, compatible con la API S3.
\end{itemize}

% ==========================================================
% 2. ARQUITECTURA Y TECNOLOGÍAS
% ==========================================================
\section{Arquitectura del sistema y tecnologías utilizadas}

La solución utiliza una arquitectura \textbf{cliente--servidor desacoplada con despliegue híbrido}. El frontend está publicado en \textbf{Vercel}; el backend y sus servicios de infraestructura se ejecutan en un \textbf{servidor local de producción/homelab} administrado por el equipo mediante Docker.

\subsection{Tecnologías principales}

\begin{table}[H]
\centering
\small
\rowcolors{2}{softgray}{white}
\begin{tabularx}{\textwidth}{@{}p{0.27\textwidth} X@{}}
\toprule
\textbf{Componente} & \textbf{Tecnología utilizada} \\
\midrule
Frontend & HTML5, CSS3, Bootstrap 4.5.3, AdminLTE 3.1.0-rc y JavaScript ES6 modular \\
Mapas & MapLibre GL JS 5.24.0 con OpenFreeMap \\
Backend & PHP 8.4+ y Laravel 13.16+ \\
Autenticación & Laravel Sanctum 4.3.2+, Bearer Tokens, verificación de correo, Google Auth y 2FA \\
Tiempo real & Laravel Reverb 1.10+ y WebSockets \\
Base de datos & PostgreSQL 16 con PostGIS 3.4 \\
Caché, sesiones y colas & Redis 7 \\
Almacenamiento & RustFS compatible con S3 \\
Servidor web & Nginx 1.27 Alpine \\
Contenedores & Docker y Docker Compose \\
Pruebas y calidad & PHPUnit 12, Vitest 3, k6 y SonarQube \\
\bottomrule
\end{tabularx}
\caption{Stack tecnológico correspondiente a la versión documentada.}
\end{table}

\subsection{Esquema general de arquitectura}

El usuario accede al frontend desde Vercel. La interfaz envía solicitudes HTTPS hacia Nginx en el servidor local de producción, que funciona como proxy reverso para la API Laravel. La aplicación utiliza PostgreSQL/PostGIS para persistencia, Redis para caché, sesiones y colas, RustFS para adjuntos y Reverb para comunicación en tiempo real.

\begin{figure}[H]
    \centering
    \includegraphics[width=0.96\textwidth,keepaspectratio]{grafico.png}
    \caption{Arquitectura general del Sistema Web de Gestión de Incidencias Georreferenciadas.}
    \label{fig:arquitectura-general}
\end{figure}

% ==========================================================
% 3. LÓGICA DE NEGOCIO
% ==========================================================
\section{Lógica de negocio por módulo}

Esta sección describe el comportamiento funcional real del sistema desde la creación
de una cuenta hasta los módulos de analítica y auditoría. Cada flujo está
implementado como un Caso de Uso dentro de la capa de Aplicación y se apoya
en entidades y reglas de dominio que son independientes del framework.

% ----------------------------------------------------------
\subsection{Módulo 1: Registro y autenticación de usuarios}
% ----------------------------------------------------------

El sistema ofrece dos mecanismos de acceso: formulario de correo y contraseña, y
autenticación con Google. Ambos comparten las mismas reglas de seguridad.

\subsubsection*{Registro con correo y contraseña}

\begin{enumerate}
    \item El usuario completa el formulario con nombre, apellido, nombre de usuario,
          correo y contraseña.
    \item El caso de uso \texttt{RegisterUserUseCase} crea la cuenta con estado
          activo pero sin correo verificado, asigna automáticamente el rol
          \textbf{CIUDADANO} y registra una identidad de tipo \texttt{local}
          para trazabilidad del proveedor de acceso.
    \item Se envía un correo de verificación de forma inmediata. Si el envío falla,
          el error se registra en el log y la cuenta queda creada igualmente; el
          usuario puede solicitar el reenvío desde el cliente.
    \item \textbf{El inicio de sesión queda bloqueado hasta que el correo sea verificado.}
          Al hacer clic en el enlace de verificación, el caso de uso
          \texttt{VerifyEmailUseCase} marca el campo \texttt{email\_verified\_at}.
\end{enumerate}

\subsubsection*{Inicio de sesión con correo y contraseña}

El caso de uso \texttt{LoginUseCase} aplica las siguientes validaciones en orden:
\begin{enumerate}
    \item Busca el usuario por correo. Si no existe, registra un intento fallido y
          lanza una excepción de credenciales inválidas (el mensaje no distingue
          si fue el correo o la contraseña, para evitar enumeración de usuarios).
    \item Verifica la contraseña mediante el \texttt{PasswordHasherPort}.
    \item Comprueba que la cuenta esté activa.
    \item Comprueba que el correo esté verificado.
    \item Si todas las comprobaciones pasan, registra un intento exitoso, actualiza
          el campo \texttt{last\_login} y sincroniza la identidad local.
    \item Si el usuario tiene \textbf{2FA habilitado}, el caso de uso no emite un
          token de sesión, sino un \texttt{two\_factor\_token} temporal. El cliente
          debe usarlo para completar el segundo factor antes de obtener acceso.
    \item Sin 2FA, emite el Bearer token de Sanctum y devuelve el perfil completo
          del usuario (roles, permisos, menús de navegación).
\end{enumerate}

Todo intento de inicio de sesión --- exitoso o fallido, con cualquier razón --- queda
registrado en la tabla \texttt{auth.login\_attempts} con correo, IP, user-agent,
resultado y motivo del fallo.

\subsubsection*{Autenticación con Google}

El caso de uso \texttt{GoogleRegistrationUseCase} acepta un \texttt{id\_token}
de Firebase y aplica la siguiente lógica:
\begin{enumerate}
    \item Verifica el token con la interfaz \texttt{GoogleTokenVerifierInterface},
          cuya implementación concreta usa la SDK de Firebase Admin.
    \item Extrae el correo, el UID de Firebase y el campo \texttt{email\_verified}
          del payload. Si alguno es inválido, el intento se registra y se rechaza.
    \item Distingue entre intención \textbf{register} e intención \textbf{login}.
          Si el usuario ya existe y se intenta registrar, o no existe y se intenta
          iniciar sesión, lanza la excepción correspondiente.
    \item En registro nuevo: crea la cuenta con contraseña aleatoria (el usuario
          nunca la conoce), asigna rol CIUDADANO y, si Google provee una foto de
          perfil, la descarga de forma asíncrona mediante el job
          \texttt{DownloadGoogleProfilePhotoJob}.
    \item En ambos flujos, sincroniza la identidad de proveedor \texttt{google}
          con el UID de Firebase, marca el correo como verificado si aún no lo está
          y actualiza el último acceso.
    \item Si el usuario tiene 2FA, aplica el mismo desvío de token temporal
          que en el login por contraseña.
\end{enumerate}

\subsubsection*{Recuperación de contraseña}

El flujo tiene tres pasos secuenciales obligatorios:
\begin{enumerate}
    \item \texttt{RequestPasswordResetCodeUseCase}: genera un código numérico de
          6 dígitos y lo envía al correo. El código tiene una ventana de validez
          definida en configuración.
    \item \texttt{VerifyPasswordResetCodeUseCase}: valida el código e intenta
          consumir el token, devolviendo una confirmación de que el código es
          correcto sin aún cambiar la contraseña.
    \item \texttt{ResetPasswordWithCodeUseCase}: con el código validado, aplica
          la nueva contraseña. Invalida el token para que no pueda reutilizarse.
\end{enumerate}

\subsubsection*{Autenticación de dos factores (2FA)}

\begin{itemize}
    \item \textbf{Habilitar:} \texttt{EnableTwoFactorUseCase} genera un secreto
          TOTP y devuelve la URL del código QR para escanear con una app
          autenticadora (Google Authenticator, etc.).
    \item \textbf{Confirmar:} \texttt{ConfirmTwoFactorUseCase} valida el primer
          código TOTP ingresado por el usuario. Solo tras esta confirmación el 2FA
          queda activo en la cuenta.
    \item \textbf{Verificar en login:} \texttt{VerifyTwoFactorLoginUseCase} recibe
          el \texttt{two\_factor\_token} temporal y el código TOTP de 6 dígitos,
          los valida y, si son correctos, emite el Bearer token definitivo.
    \item \textbf{Deshabilitar:} \texttt{DisableTwoFactorUseCase} elimina el
          secreto almacenado y desactiva el flag en la cuenta.
\end{itemize}

% ----------------------------------------------------------
\subsection{Módulo 2: Gestión de usuarios y perfiles}
% ----------------------------------------------------------

El administrador gestiona la base de usuarios del sistema. Un usuario tiene
nombre, apellido, nombre de usuario, correo, estado activo/inactivo y rol.
Los datos se validan mediante \texttt{FormRequest} antes de llegar al caso de uso.

Cada usuario autenticado puede actualizar su propio perfil
(\texttt{UpdateOwnProfileUseCase}), cambiar su contraseña
(\texttt{ChangeOwnPasswordUseCase}) y gestionar su foto de perfil
(\texttt{UpdateOwnProfilePhotoUseCase} y \texttt{GetOwnProfilePhotoUseCase}).

Los roles disponibles son: \textbf{ADMIN}, \textbf{SUPERVISOR}, \textbf{OPERADOR}
y \textbf{CIUDADANO}. Cada rol tiene un conjunto de permisos granulares asignados
que determinan qué operaciones puede ejecutar el usuario en cada endpoint. El
sistema evalúa permisos --- no solo roles --- antes de procesar cualquier acción
sensible.

% ----------------------------------------------------------
\subsection{Módulo 3: Estructura operativa}
% ----------------------------------------------------------

La estructura operativa relaciona supervisores, operadores y territorios. Esto
define qué personas atienden incidencias en qué zonas geográficas.

\begin{itemize}
    \item \textbf{Zonas operativas:} cada zona abarca una o más unidades
          territoriales. Un supervisor puede ser asignado a una zona
          (\texttt{assignSupervisorToZone}).
    \item \textbf{Supervisores y operadores:} un supervisor tiene un límite
          configurable de operadores a cargo, controlado por el servicio de dominio
          \texttt{SupervisorCapacityPolicy}. Si se intenta asignar más operadores
          que el límite configurado, o si se reduce el límite por debajo del número
          actual de operadores, el dominio lanza una excepción de negocio.
    \item \textbf{Territorios:} a cada operador se le asigna un territorio de
          trabajo (\texttt{assignOperatorTerritory}). El sistema filtra la
          visibilidad de incidencias según el territorio del operador.
    \item \textbf{Perfiles:} tanto supervisores como operadores tienen un perfil
          con datos operativos que pueden actualizarse de forma independiente al
          perfil de autenticación del usuario.
\end{itemize}

% ----------------------------------------------------------
\subsection{Módulo 4: Registro de incidencias}
% ----------------------------------------------------------

Una incidencia es el objeto central del sistema. Puede ser reportada por cualquier
usuario con acceso a la plataforma.

\subsubsection*{Creación}

El caso de uso \texttt{IncidentUseCase::store} persiste la incidencia con los
datos básicos: título, descripción, coordenadas geográficas (latitud y longitud)
y unidad territorial. En la base de datos, el trigger
\texttt{trg\_actualizar\_ubicacion} sincroniza automáticamente el campo
\texttt{geometry(Point, 4326)} de PostGIS con las coordenadas proporcionadas,
garantizando la consistencia entre los campos numéricos y la representación
espacial.

La incidencia se crea con estado inicial \textbf{NUEVA} y con
\texttt{classification\_status = Pending}, lo que indica que aún no ha sido
clasificada por un supervisor o administrador.

\subsubsection*{Campos requeridos y opcionales}

\begin{itemize}
    \item \textbf{Obligatorios:} título, descripción y categoría.
    \item \textbf{Opcionales en creación:} subcategoría, prioridad, dirección,
          latitud, longitud, unidad territorial, fecha de resolución y detalle
          de clasificación.
\end{itemize}

\subsubsection*{Adjuntos}

El método \texttt{attachFile} recibe el archivo validado por la capa HTTP,
lo sube al almacenamiento S3 compatible (RustFS) mediante el puerto
\texttt{FileStoragePort} y registra únicamente la metadata en la base de datos:
nombre original, ruta en el bucket, tipo MIME, tamaño en bytes y hash de
integridad. Los archivos binarios nunca se guardan en PostgreSQL.

\subsubsection*{Comentarios}

Cualquier usuario con acceso a la incidencia puede añadir comentarios con texto
y fecha mediante \texttt{addComment}. Cada comentario queda ligado al usuario
que lo creó y a la incidencia correspondiente.

\subsubsection*{Edición y eliminación}

La edición solo está permitida cuando la entidad de dominio \texttt{Incident}
confirma que el estado actual lo permite (\texttt{canBeEdited()}). Si el estado
no lo permite, el caso de uso lanza \texttt{IncidentException::editNotAllowed}
sin llegar a la capa de persistencia. La eliminación es lógica (soft delete).

% ----------------------------------------------------------
\subsection{Módulo 5: Clasificación de incidencias}
% ----------------------------------------------------------

Después de que un ciudadano reporta una incidencia, un supervisor o administrador
debe clasificarla. La clasificación asigna categoría, subcategoría, prioridad
y cualquier otro metadato necesario para la atención.

El método \texttt{classify} del caso de uso registra el historial de clasificación
en \texttt{incident\_classification\_history}, de modo que queda trazabilidad de
todos los cambios de categoría y subcategoría, con la razón del cambio, el usuario
responsable y la fecha.

Hasta que la incidencia sea clasificada, \texttt{classification\_status} permanece
en \texttt{Pending} y la entidad de dominio rechaza las operaciones que requieren
clasificación previa (\texttt{requiresClassification()} retorna \texttt{true}).

\subsubsection*{Solicitud de nueva categoría}

Si al momento de clasificar una incidencia el supervisor no encuentra una categoría
apropiada en el catálogo existente, tiene la facultad de proponer la creación de
una nueva categoría al equipo administrador.

\begin{enumerate}
    \item El supervisor envía una solicitud utilizando \texttt{CategoryRequestUseCase::requestNewCategory},
          proporcionando un nombre sugerido para la categoría, una justificación de la necesidad y,
          opcionalmente, vinculando la incidencia actual.
    \item La solicitud se registra con estado \textbf{NUEVA} (\texttt{pending}) y
          se notifica a todos los administradores activos con permiso para gestionar catálogos,
          utilizando el servicio \texttt{AdminNotifier}.
    \item El equipo administrador revisa la bandeja de solicitudes pendientes. Un administrador
          puede:
          \begin{itemize}
              \item \textbf{Aprobar:} El caso de uso \texttt{CategoryRequestUseCase::approveRequest}
                    actualiza el estado de la solicitud a aprobada y registra al administrador que
                    tomó la decisión. Tras esto, el administrador debe crear manualmente la categoría
                    en el catálogo. El supervisor solicitante recibe una notificación de aprobación.
              \item \textbf{Rechazar:} El caso de uso \texttt{CategoryRequestUseCase::rejectRequest}
                    actualiza el estado a rechazada, exigiendo un comentario o justificación. El supervisor
                    solicitante es notificado del rechazo y del motivo correspondiente.
          \end{itemize}
\end{enumerate}

% ----------------------------------------------------------
\subsection{Módulo 6: Estados y ciclo de vida de la incidencia}
% ----------------------------------------------------------

El ciclo de vida de una incidencia está gobernado por un grafo de estados
persistido en la tabla \texttt{state\_transitions}. El sistema no tiene
transiciones hardcodeadas: las rutas válidas entre estados se configuran en
la base de datos y se consultan en tiempo de ejecución.

\subsubsection*{Estados del sistema}

\begin{table}[H]
\centering
\small
\rowcolors{2}{softgray}{white}
\begin{tabularx}{\textwidth}{@{}p{0.22\textwidth} X@{}}
\toprule
\textbf{Estado} & \textbf{Significado operativo} \\
\midrule
NUEVA & Incidencia recién reportada, sin atención aún. \\
ASIGNADA & Se han designado operadores responsables. \\
EN PROGRESO & Un operador está trabajando activamente en la atención. \\
RESUELTA & El operador solicitó resolución y fue aprobada por el supervisor. \\
CERRADA & Ciclo finalizado sin posibilidad de reapertura inmediata. \\
REABIERTA & El supervisor o administrador reabrió la incidencia; inicia un nuevo ciclo. \\
RECHAZADA & La incidencia fue evaluada y no procede su atención. \\
CANCELADA & La incidencia fue cancelada antes de ser atendida. \\
\bottomrule
\end{tabularx}
\caption{Estados de incidencia y su significado operativo.}
\end{table}

\subsubsection*{Cambio de estado directo}

El método \texttt{changeState} ejecuta las siguientes verificaciones antes de
persistir cualquier cambio:
\begin{enumerate}
    \item Carga la incidencia con bloqueo de escritura (\texttt{SELECT FOR UPDATE})
          para evitar condiciones de carrera.
    \item Busca la transición entre el estado actual y el estado destino en el
          repositorio. Si no existe, lanza \texttt{IncidentException::transitionNotAllowed}.
    \item Verifica que los roles del usuario autenticado estén incluidos en los
          roles permitidos por esa transición. Si no, lanza
          \texttt{IncidentException::transitionForbidden}.
    \item Si el estado destino es REABIERTA, verifica que el usuario tenga el
          permiso \texttt{incidents.reopen} además de su rol.
    \item Si la transición requiere comentario obligatorio y el comentario está
          vacío o es solo espacios en blanco, lanza
          \texttt{IncidentException::transitionRequiresComment}.
    \item Aplica el cambio de estado dentro de una transacción.
\end{enumerate}

\subsubsection*{Solicitud de resolución por operador}

El operador no puede marcar una incidencia como resuelta directamente. En su lugar,
utiliza el flujo de solicitud de cambio de estado:
\begin{enumerate}
    \item El operador llama a \texttt{requestStateChange}, que valida que el
          solicitante tenga rol OPERADOR, que la incidencia esté en un estado que
          permita solicitar resolución (\texttt{canRequestResolution()}), que el
          estado destino sea RESUELTA, que el operador tenga una asignación activa
          en esa incidencia, y que no exista ya una solicitud pendiente.
    \item Se crea la solicitud y se notifica al supervisor responsable en tiempo
          real mediante WebSockets (Reverb).
    \item El supervisor puede \textbf{aprobar} (\texttt{approveStateChange}) o
          \textbf{rechazar} (\texttt{rejectStateChange}) la solicitud. En ambos
          casos, el operador recibe una notificación con el resultado.
    \item Solo los roles SUPERVISOR y ADMIN pueden aprobar o rechazar solicitudes.
\end{enumerate}

\subsubsection*{Reapertura y ciclos}

Cuando una incidencia es reabierta, el sistema crea un nuevo registro en
\texttt{incident\_cycles}. Las asignaciones del ciclo anterior se desactivan y se
clonan al nuevo ciclo, preservando la trazabilidad histórica completa. Esto
permite saber qué operadores atendieron cada ciclo de vida de la incidencia.

% ----------------------------------------------------------
\subsection{Módulo 7: Asignación de operadores}
% ----------------------------------------------------------

El supervisor (o administrador) asigna operadores a una incidencia mediante
\texttt{assign}. La asignación registra un operador responsable principal y
opcionalmente operadores de apoyo. Cada asignación queda vinculada al ciclo
de vida activo de la incidencia.

Para facilitar la selección, el método \texttt{assignmentOperatorOptions} devuelve
la lista de operadores disponibles según la visibilidad del supervisor. El sistema
distingue entre asignación activa e inactiva: al reabrir una incidencia, las
asignaciones del ciclo anterior se inactivan automáticamente.

Una vez que un operador tiene una asignación activa, puede solicitar el cambio de
estado hacia RESUELTA. Si no tiene asignación activa, la solicitud es rechazada
por el dominio antes de llegar a la base de datos.

% ----------------------------------------------------------
\subsection{Módulo 8: Notificaciones en tiempo real}
% ----------------------------------------------------------

El sistema genera notificaciones automáticas ante eventos relevantes del ciclo de
vida de una incidencia. Las notificaciones se almacenan en la tabla
\texttt{notifications} con referencia al usuario destinatario, el texto del mensaje
y el estado de lectura.

Los eventos que generan notificaciones incluyen: solicitud de cambio de estado,
aprobación de solicitud, rechazo de solicitud, y cualquier otro evento configurado
en el puerto \texttt{IncidentStateChangeNotifierPort}.

La entrega en tiempo real se realiza mediante \textbf{Laravel Reverb} (WebSockets).
El frontend recibe el evento instantáneamente y actualiza el contador de no leídas
sin recargar la página.

El usuario puede:
\begin{itemize}
    \item Consultar todas sus notificaciones paginadas.
    \item Ver el conteo de no leídas en tiempo real.
    \item Marcar una notificación individual como leída (\texttt{markAsRead}).
    \item Marcar todas las notificaciones como leídas de una vez
          (\texttt{markAllAsRead}).
\end{itemize}

% ----------------------------------------------------------
\subsection{Módulo 9: Mapa georreferenciado}
% ----------------------------------------------------------

El método \texttt{mapPoints} devuelve la lista de incidencias con sus coordenadas
geográficas para su representación en el mapa interactivo. La consulta aplica
el mismo filtro de visibilidad por rol y territorio que el listado de tabla.

El backend expone \texttt{GET /api/incidents/map} con la colección de puntos.
El frontend los renderiza mediante \textbf{MapLibre GL JS} sobre teselas de
\textbf{OpenFreeMap} (sin token, sin costo). El ciudadano puede usar el mapa
directamente al reportar una incidencia: hace clic en el punto del mapa, las
coordenadas se capturan automáticamente y se completan los campos de latitud y
longitud en el formulario.

% ----------------------------------------------------------
\subsection{Módulo 10: Catálogos y unidades territoriales}
% ----------------------------------------------------------

Los catálogos son datos de referencia que el sistema consume desde la API;
nada está codificado directamente en la interfaz.

\begin{itemize}
    \item \textbf{Categorías y subcategorías:} clasifican el tipo de incidencia.
    \item \textbf{Prioridades:} cada prioridad tiene un nombre y un SLA
          (en horas) que determina el tiempo esperado de resolución.
    \item \textbf{Estados y transiciones:} se consultan para construir dinámicamente
          los controles de cambio de estado disponibles para el usuario.
    \item \textbf{Roles y permisos:} se leen de la API para construir la interfaz
          de administración.
    \item \textbf{Unidades territoriales:} organizadas en tres niveles jerárquicos
          --- Provincia, Cantón, Parroquia --- con código de referencia y estado
          activo/inactivo. El administrador puede crear, editar y desactivar
          unidades. La interfaz carga las unidades en cascada (provincias →
          cantones → parroquias) según la selección del usuario.
\end{itemize}

% ----------------------------------------------------------
\subsection{Módulo 11: Dashboard y analítica}
% ----------------------------------------------------------

El dashboard presenta indicadores clave de rendimiento (KPI) obtenidos mediante
\texttt{GetDashboardMetricsUseCase}. Los reportes avanzados provienen de
\texttt{ReportUseCase::getAnalytics}, que procesa las incidencias aplicando el
filtro de visibilidad del usuario y calcula:

\begin{itemize}
    \item Total de incidencias, tasa de resolución y tiempo promedio de resolución
          en días.
    \item Incidencias activas, resueltas, cerradas, vencidas, críticas y de alta
          prioridad.
    \item Incidencias reportadas en los últimos 7 días.
    \item Distribución por prioridad y por categoría.
    \item Top 5 de unidades territoriales con más incidencias.
    \item Tendencia mensual de incidencias registradas, resueltas y pendientes
          (agrupadas por mes).
    \item Tabla resumen por categoría con totales, pendientes, resueltos y tasa
          de resolución.
    \item Tiempo promedio de resolución por categoría.
\end{itemize}

Los reportes se pueden filtrar por rango de fechas, categoría y estado. La
visibilidad de los datos respeta el rol: un operador solo ve las incidencias
de su territorio, mientras que supervisores y administradores tienen visión global.

% ----------------------------------------------------------
\subsection{Módulo 12: Auditoría y trazabilidad del sistema}
% ----------------------------------------------------------

El módulo de auditoría expone dos consultas paginadas accesibles solo para roles
con el permiso correspondiente:

\begin{itemize}
    \item \textbf{Logs de actividad:} registran operaciones de crear, actualizar
          y eliminar sobre las tablas del sistema, con referencia al usuario, la
          tabla afectada, el ID del registro y la acción realizada.
    \item \textbf{Intentos de autenticación:} registran todos los intentos de
          login --- exitosos y fallidos --- con correo, IP, user-agent, resultado
          y, en caso de fallo, el motivo (usuario no encontrado, contraseña
          incorrecta, cuenta inactiva, correo no verificado, token de Google
          inválido, etc.).
\end{itemize}

Ambas consultas admiten filtros y paginación. La tabla de intentos de autenticación
es un instrumento de seguridad operacional que permite al administrador detectar
patrones de acceso no autorizados.

\textbf{Limitaciones actuales:} el catálogo territorial está especializado para
Ecuador y no modela el país como nivel independiente. El despliegue productivo
utiliza múltiples servicios Docker, pero no configura réplicas del backend para
escalamiento horizontal.


% ==========================================================
% 4. BASE DE DATOS
% ==========================================================
\section{Base de datos}

La persistencia se implementó en \textbf{PostgreSQL 16 con PostGIS}. El diseño
relacional utiliza claves primarias y foráneas, restricciones de integridad,
índices y componentes SQL para mantener la consistencia de la información.
Para organizar las responsabilidades, el modelo se distribuye principalmente
en los esquemas \texttt{core} y \texttt{auth}.

El esquema \texttt{core} concentra las entidades relacionadas con incidencias,
clasificación, estados, asignaciones, comentarios, adjuntos, notificaciones,
ciclos y unidades territoriales. El esquema \texttt{auth} concentra usuarios,
roles, permisos, sesiones, perfiles operativos e identidades. Ambos esquemas
forman parte de la misma base de datos y se relacionan durante la ejecución.

\subsection{Modelo lógico}

El modelo lógico representa las entidades principales y sus relaciones sin
depender de tipos físicos específicos de PostgreSQL. Las entidades centrales son
\texttt{Usuario} e \texttt{Incidente}: un usuario puede reportar, clasificar,
comentar o atender incidencias según sus roles y permisos; cada incidencia se
relaciona con categoría, subcategoría, prioridad, estado, unidad territorial,
asignaciones, comentarios, adjuntos, notificaciones e históricos.

El código fuente (Mermaid) de este diagrama se encuentra disponible en el repositorio bajo el archivo \texttt{DATABASE\_LOGICAL\_MODEL.mmd}.

\begin{figure}[H]
    \centering
    \includegraphics[width=0.97\textwidth,keepaspectratio]{modelo_logico.png}
    \caption{Modelo lógico general del Sistema Web de Gestión de Incidencias Georreferenciadas.}
    \label{fig:modelo-logico}
\end{figure}

\subsection{Modelo físico del esquema \texttt{core}}

El modelo físico del esquema \texttt{core} muestra la implementación real de
las tablas, columnas y relaciones utilizadas para el ciclo de vida de las
incidencias. Entre las tablas principales se encuentran
\texttt{incidents}, \texttt{states}, \texttt{state\_transitions},
\texttt{incident\_states}, \texttt{incident\_assignments},
\texttt{incident\_comments}, \texttt{incident\_attachments},
\texttt{categories}, \texttt{subcategories}, \texttt{priorities},
\texttt{territorial\_units}, \texttt{incident\_cycles},
\texttt{notifications}, \texttt{incident\_classification\_history},
\texttt{category\_requests}, \texttt{state\_change\_requests} y
\texttt{settings}.

\clearpage
\begin{landscape}
\begin{figure}[H]
    \centering
    \includegraphics[width=0.96\linewidth,height=0.82\textheight,keepaspectratio]{modelo_core.png}
    \caption{Modelo físico de la base de datos correspondiente al esquema \texttt{core}.}
    \label{fig:modelo-core}
\end{figure}
\end{landscape}
\clearpage

\subsection{Modelo físico del esquema \texttt{auth}}

El esquema \texttt{auth} contiene la implementación física de las entidades de
identidad, autenticación, autorización y estructura operativa. Incluye
\texttt{users}, \texttt{roles}, \texttt{permissions}, \texttt{role\_user},
\texttt{permission\_role}, \texttt{operator\_profiles},
\texttt{supervisor\_profiles}, \texttt{supervisor\_operator\_assignments},
\texttt{user\_territories}, \texttt{sessions}, \texttt{user\_identities},
\texttt{navigation\_items} (con soporte de \texttt{allowed\_roles} por ítem),
\texttt{password\_reset\_tokens} y \texttt{personal\_access\_tokens}.

\clearpage
\begin{landscape}
\begin{figure}[H]
    \centering
    \includegraphics[width=0.96\linewidth,height=0.82\textheight,keepaspectratio]{modelo_auth.png}
    \caption{Modelo físico de la base de datos correspondiente al esquema \texttt{auth}.}
    \label{fig:modelo-auth}
\end{figure}
\end{landscape}
\clearpage

\subsection{Relación entre los esquemas \texttt{core} y \texttt{auth}}

Aunque los esquemas se encuentran separados lógicamente, ambos colaboran
durante la ejecución. \texttt{auth} proporciona identidad, roles, permisos,
perfiles y territorios; \texttt{core} utiliza dichas referencias para registrar
quién reporta, clasifica, comenta, atiende o cambia el estado de una incidencia.

\begin{center}
\begin{tcolorbox}[
    colback=softgray,
    colframe=primarygray,
    width=0.90\textwidth,
    boxrule=0.7pt
]
\centering
\textbf{PostgreSQL 16 + PostGIS}\\[0.35cm]
\texttt{auth} $\longrightarrow$ \textbf{Usuarios, roles, permisos y estructura operativa}\\[0.20cm]
$\updownarrow$\\[0.20cm]
\texttt{core} $\longrightarrow$ \textbf{Incidencias, estados, asignaciones y seguimiento}
\end{tcolorbox}
\end{center}

\subsection{Normalización, integridad y programación SQL}

El diseño aplica normalización mediante entidades independientes para estados,
prioridades, categorías, subcategorías, usuarios, roles, permisos y unidades
territoriales, evitando duplicación innecesaria de datos. La integridad se
refuerza mediante claves primarias, claves foráneas, restricciones
\texttt{CHECK}, índices convencionales, índices parciales e índices espaciales.

PostGIS se utiliza para almacenar la ubicación geográfica mediante
\texttt{GEOMETRY(Point,4326)}. Además, la base de datos incorpora funciones y
triggers para automatizar operaciones, como la sincronización de latitud y
longitud con la geometría espacial (el cálculo automático de fecha límite por SLA
fue deshabilitado por requerimiento del flujo operativo).

\subsection{Archivo SQL de entrega}

El archivo SQL acompaña la entrega académica y se encuentra disponible directamente en la raíz del repositorio (\texttt{backup.sql}), de modo que la estructura y los datos necesarios puedan reconstruirse sin depender únicamente de las migraciones del framework.

\begin{center}
\textbf{Archivo incluido en el repositorio:} \texttt{backup.sql}
\end{center}

También puede regenerarse desde el servidor de base de datos mediante:

\begin{verbatim}
pg_dump -U user_im -d incident_management_system \
  --format=plain --no-owner --no-privileges > backup.sql
\end{verbatim}

% ==========================================================
% 5. CREDENCIALES
% ==========================================================
\section{Credenciales de acceso}

Los \textit{seeders} de demostración incluyen los siguientes usuarios, todos con contraseña \texttt{password}:

\begin{table}[H]
\centering
\small
\rowcolors{2}{softgray}{white}
\begin{tabularx}{\textwidth}{@{}l X l@{}}
\toprule
\textbf{Rol} & \textbf{Correo} & \textbf{Contraseña} \\
\midrule
Administrador & admin@incidents.local & password \\
Administrador (Revisión) & admindocente@incidents.local & password \\
Supervisor (Guayas Z2) & supervisor.guayas@incidents.local & password \\
Operador (Guayas Z2) & operador.z2.01@incidents.local & password \\
Ciudadano & ciudadano1@incidents.local & password \\
\bottomrule
\end{tabularx}
\caption{Credenciales de demostración cargadas mediante seeders.}
\end{table}

\textbf{Nota:} estas credenciales son únicamente para pruebas y demostración académica.

% ==========================================================
% 6. EJECUCIÓN
% ==========================================================
\section{Instrucciones de ejecución}

\textbf{Requisitos:} Git, Docker Desktop/Engine con Docker Compose, PHP 8.3+ y Composer.

\begin{enumerate}
    \item \textbf{Clonar el repositorio:}
\begin{verbatim}
git clone https://github.com/DevTorres404/georeferenced-incident-management.git
cd georeferenced-incident-management
\end{verbatim}

    \item \textbf{Levantar servicios de apoyo del backend:}
\begin{verbatim}
cd backend
docker compose -f docker-compose.yml up -d
\end{verbatim}
    Este entorno inicia PostgreSQL/PostGIS, Redis, RustFS y Mailpit.

    \item \textbf{Configurar Laravel:}
\begin{verbatim}
composer install
cp .env.example .env
php artisan key:generate
php artisan migrate:fresh --seed
\end{verbatim}

    \item \textbf{Iniciar backend y servicios de ejecución:}
\begin{verbatim}
php artisan serve
php artisan queue:listen
php artisan reverb:start
\end{verbatim}

    \item \textbf{Levantar el frontend local con Nginx:}
\begin{verbatim}
cd ../frontend
docker compose up -d
\end{verbatim}

\end{enumerate}



% ==========================================================
% 7. DESPLIEGUE
% ==========================================================
\section{Despliegue}

El sistema utiliza un \textbf{despliegue híbrido}. El frontend se encuentra
publicado en \textbf{Vercel}, mientras que el backend y sus servicios de
infraestructura se ejecutan en un \textbf{servidor local de producción/homelab}
administrado por el equipo.

El entorno productivo del backend utiliza el archivo
\texttt{docker-compose.prod.yml}. En él se definen los servicios
\texttt{app}, \texttt{migrate}, \texttt{nginx}, \texttt{queue},
\texttt{scheduler}, \texttt{reverb}, \texttt{pgsql}, \texttt{redis},
\texttt{rustfs} y \texttt{rustfs-init}. PostgreSQL/PostGIS y Redis disponen de
\textit{healthchecks}; los datos de PostgreSQL, Redis, almacenamiento y
directorios de la aplicación se conservan mediante volúmenes Docker.

Nginx funciona como punto de entrada HTTP del backend. PostgreSQL, Redis y
RustFS permanecen dentro de la red \textit{bridge} de Docker y no requieren
exposición directa al usuario final. Los trabajadores de cola procesan tareas
asíncronas, el scheduler ejecuta tareas programadas y Reverb mantiene la
comunicación WebSocket.

\textbf{URL funcional del sistema:}
\begin{center}
\url{https://app.labtorres.me/}
\end{center}

\clearpage
\begin{figure}[H]
    \centering
    \includegraphics[width=0.88\textwidth,height=0.34\textheight,keepaspectratio]{vercel_prod.png}
    \caption{Frontend del sistema desplegado y disponible en Vercel.}
    \label{fig:vercel-prod}
\end{figure}

\vspace{0.15cm}

\begin{figure}[H]
    \centering
    \includegraphics[width=0.96\textwidth,height=0.34\textheight,keepaspectratio]{docker_prod.png}
    \caption{Servicios Docker del backend ejecutándose en el servidor local de producción.}
    \label{fig:docker-prod}
\end{figure}



% ==========================================================
% 8. EVIDENCIAS Y CALIDAD
% ==========================================================
\section{Evidencias de implementación y calidad del sistema}

Esta sección reúne las evidencias solicitadas para demostrar funcionamiento,
validación, testing, rendimiento y análisis de calidad. Las capturas se mantienen
como figuras independientes, pero se organizan de dos por página cuando su
resolución permite una lectura adecuada.

\subsection{Capturas funcionales del sistema}

Las evidencias funcionales cubren la pantalla principal, el registro y gestión
de incidencias, la consulta georreferenciada y el módulo de analítica.

\dosporevidencia
{dashboard.png}{0.91}{Dashboard principal con indicadores y visualización de información.}
{registro.png}{0.91}{Formulario de registro de una incidencia con selección geográfica.}

\dosporevidencia
{detalle.png}{0.91}{Detalle de incidencia con seguimiento, estado, asignaciones e historial.}
{mapa.png}{0.94}{Mapa y consulta de incidencias georreferenciadas mediante MapLibre GL JS.}

\clearpage
\begin{figure}[H]
    \centering
    \includegraphics[width=0.93\textwidth,height=0.62\textheight,keepaspectratio]{reportes.png}
    \caption{Vista de analítica, filtros, métricas y reportes del sistema.}
\end{figure}

\subsection{Validaciones y manejo de errores}

La validación se aplica en diferentes niveles:

\begin{itemize}
    \item \textbf{Frontend:} validación de datos antes del envío, estados de
    carga y mensajes de error orientados al usuario.
    \item \textbf{Backend:} reglas de validación de Laravel, autenticación,
    autorización basada en permisos y manejo de operaciones no permitidas.
    \item \textbf{Dominio:} reglas para transiciones de estado, clasificación,
    asignaciones, resolución y reapertura.
    \item \textbf{Base de datos:} claves foráneas, restricciones, índices,
    funciones y triggers que refuerzan la integridad.
\end{itemize}

\subsection{Casos de prueba y pruebas automatizadas}

La versión documentada incluye pruebas automatizadas tanto en backend como en
frontend. La evidencia de PHPUnit muestra \textbf{220 pruebas superadas},
\textbf{1.357 aserciones} y \textbf{0 fallos}. La ejecución registrada tuvo una
duración de \textbf{501,41 s}.

La evidencia de Vitest corresponde al comando
\texttt{npm run test:js:coverage}. Se observan \textbf{54 archivos de prueba
superados}, \textbf{1.275 tests superados} y \textbf{4 omitidos}, para un total
de 1.279 tests registrados.

\dosporevidencia
{phpunit.png}{0.96}{Resultado de la suite PHPUnit del backend: 220 pruebas superadas y 1.357 aserciones.}
{vitest.png}{0.96}{Resultado de las pruebas del frontend con Vitest: 51 archivos superados y 1.275 tests aprobados.}

\subsection{Prueba de carga o estrés}

El repositorio incluye escenarios k6 para salud del servicio, autenticación,
mapa y flujo completo: \texttt{e1-health.js}, \texttt{e2-login.js},
\texttt{e3-map.js} y \texttt{e4-full-flow.js}.

La evidencia incluida corresponde al escenario \texttt{e1-health.js}. La
ejecución alcanzó hasta \textbf{50 VUs}, realizó \textbf{3.447 solicitudes HTTP}
y \textbf{6.894 checks}. Se obtuvo \textbf{100\% de checks exitosos},
\textbf{0\% de solicitudes fallidas}, un tiempo promedio de
\textbf{163,13 ms}, \textbf{p95 de 197,47 ms} y un máximo de
\textbf{311,2 ms}. Los umbrales configurados para tasa de error y p95 fueron
superados satisfactoriamente.

\subsection{Uso de herramientas de calidad y métricas}

Se utilizó \textbf{SonarQube} como herramienta de análisis de calidad. En la
captura final se observan los siguientes indicadores:

\begin{table}[H]
\centering
\small
\rowcolors{2}{softgray}{white}
\begin{tabularx}{\textwidth}{@{}p{0.38\textwidth} X@{}}
\toprule
\textbf{Indicador} & \textbf{Resultado observado} \\
\midrule
Security & 0 incidencias abiertas, calificación A \\
Reliability & 6 incidencias abiertas, calificación C \\
Maintainability & 352 incidencias abiertas, calificación A \\
Coverage & 74,4\% sobre aproximadamente 23k líneas a cubrir \\
Duplications & 1,3\% sobre aproximadamente 59k líneas \\
Security Hotspots & 0, calificación A \\
Accepted issues & 0 \\
\bottomrule
\end{tabularx}
\caption{Métricas observadas en el análisis final de SonarQube.}
\end{table}

\dosporevidencia
{k6.png}{0.96}{Resultado de la prueba de carga k6 del escenario de salud del servicio.}
{sonarqube.png}{0.96}{Panel de métricas de calidad del proyecto en SonarQube.}

\subsection{Síntesis de evidencias de calidad}

Las pruebas funcionales y automatizadas verifican los módulos principales; k6
proporciona evidencia de comportamiento bajo carga; SonarQube aporta métricas
de cobertura, duplicación, seguridad, fiabilidad y mantenibilidad. Las
dificultades técnicas detectadas durante el desarrollo y las soluciones
aplicadas se documentan en la Sección 9.


% ==========================================================
% 9. DIFICULTADES Y SOLUCIONES
% ==========================================================
\section{Dificultades encontradas y soluciones}

\begin{table}[H]
\centering
\small
\begin{tabularx}{\textwidth}{@{}p{0.30\textwidth} X@{}}
\toprule
\textbf{Dificultad} & \textbf{Solución aplicada} \\
\midrule
Consistencia entre coordenadas y geometría espacial & PostGIS y un trigger mantienen sincronizados latitud/longitud y la geometría \texttt{Point}. \\
Trazabilidad de cambios de estado & Se separó el estado actual del historial y se añadieron transiciones, solicitudes y ciclos de reapertura. \\
Acoplamiento entre pantallas y API & El consumo HTTP se centralizó en un cliente de API y el frontend se organizó en módulos. \\
Procesamiento asíncrono & Redis, workers de cola y Reverb permiten ejecutar tareas y notificaciones sin bloquear las peticiones principales. \\
Reproducibilidad del entorno & Docker Compose separa servicios, redes, volúmenes y configuración mediante variables de entorno. \\
\bottomrule
\end{tabularx}
\caption{Principales dificultades técnicas y soluciones aplicadas.}
\end{table}

% ==========================================================
% 10. CONCLUSIONES
% ==========================================================
\section{Conclusiones}

El proyecto integró Desarrollo Web, Base de Datos, Calidad de Software y Administración de Data Center en una solución funcional. La separación entre frontend y backend permitió desplegar cada componente de manera independiente, mientras PostgreSQL/PostGIS proporcionó persistencia relacional y geoespacial.

El historial de estados, asignaciones, comentarios, notificaciones, clasificación y ciclos de resolución permite mantener la trazabilidad del ciclo de vida de una incidencia. Los roles y permisos controlan el acceso a las operaciones y la estructura territorial permite relacionar el trabajo operativo con la ubicación de los reportes.

En calidad se evidenciaron pruebas automatizadas de backend y frontend, pruebas de carga con k6 y análisis con SonarQube. En infraestructura, Docker Compose permitió desplegar de manera reproducible la API, Nginx, PostgreSQL/PostGIS, Redis, Reverb, workers, scheduler y almacenamiento S3 compatible.

Como resultado, el SGI supera el alcance de un CRUD básico e integra seguimiento, georreferenciación, trazabilidad, seguridad, métricas, reportes y despliegue productivo. Como mejoras futuras quedan el escalamiento horizontal del backend y la generalización del catálogo territorial para soportar múltiples países.


% ==========================================================


\end{document}
